% Informe final del TP de Visión Artificial (#27). Molde con la estructura del TP2.
%
% Cada sección dice quién la escribe (reparto del #27) y de dónde salen los datos.
% Lo que falta completar está en rojo con \pend{...}: antes de entregar, buscar "\pend" y que no quede ninguno.
% Las figuras van en informe/figuras/. Si falta el logo (figuras/logo_udesa.png), la portada sale sin él.
%
% Compilar: pdflatex informe.tex (dos veces, para las referencias cruzadas), o subir la carpeta a Overleaf.

\documentclass[10pt,a4paper]{article}

\usepackage[utf8]{inputenc}
\usepackage[T1]{fontenc}
\usepackage[spanish,es-tabla,es-nodecimaldot]{babel}
\usepackage{lmodern}
\usepackage[a4paper,margin=2cm,columnsep=0.8cm]{geometry}
\usepackage{amsmath,amssymb}
\usepackage{graphicx}
\usepackage{booktabs}
\usepackage{multirow}
\usepackage{xcolor}
\usepackage[font=small,labelfont=bf]{caption}
\usepackage{enumitem}
\usepackage[hidelinks]{hyperref}

\graphicspath{{figuras/}}

% Texto pendiente: se ve en rojo hasta que se complete.
\newcommand{\pend}[1]{\textcolor{red}{[#1]}}
% Figura pendiente: un recuadro con la descripción, del ancho de la columna.
\newcommand{\figpend}[1]{\fbox{\parbox[c][3.5cm][c]{0.92\columnwidth}{\centering\small\textcolor{red}{#1}}}}

\begin{document}

% ------------------------------------------------------------------ portada

\begin{titlepage}
\centering
\IfFileExists{figuras/logo_udesa.png}{\includegraphics[width=0.35\textwidth]{logo_udesa.png}\par}{}
\vspace{1.5cm}
{\LARGE\bfseries Trabajo Práctico Final:\\[0.2em] \pend{título distintivo}\par}
\vspace{0.4cm}
{\large Alerta temprana de actividad en Vaca Muerta con imágenes satelitales\par}
\vspace{1cm}
{\Large Visión Artificial\par}
\vspace{1cm}
Grupo 4: Sofía Bottino, Agustina Barrena, Bianca Paz\par
\vspace{0.8cm}
Información de contacto: sbottino@udesa.edu.ar, abarrena@udesa.edu.ar, bpaz@udesa.edu.ar\par
\vspace{0.8cm}
Segundo semestre 2026\par
\vspace{1.5cm}

\begin{flushleft}
{\large\bfseries Resumen\par}
\vspace{0.3cm}
% Entre las tres, al final. Un párrafo como el del TP2: problema, datos, método y los números principales.
\pend{Resumen (entre las tres, se escribe al final). Problema: detectar locaciones petroleras nuevas de Vaca Muerta
en imágenes Sentinel-2 antes del inicio de perforación. Datos: N pares validados de dos zonas (Añelo y zona externa).
Método: red siamesa con encoder ResNet34 y K-fold por bloques. Resultados: IoU de locación en test\_anelo y
test\_externo, \% de locaciones detectadas a tiempo y mediana de la anticipación en meses.}
\end{flushleft}
\end{titlepage}

\twocolumn

% ------------------------------------------------------------------ 1

\section{Introducción y Marco Teórico}
% Sofi (#23). Fuentes: docs/PROBLEMA.md, ADR-0002, docs/ESTADO_DEL_ARTE.md. Vocabulario de GLOSSARY.md.

En el presente trabajo estudiamos la detección de locaciones petroleras nuevas en imágenes satelitales,
un problema de detección de cambios con aplicaciones en inteligencia de mercado y en el monitoreo de la huella
de la industria. Lo abordamos con el objetivo de anticipar la actividad de perforación en Vaca Muerta: entrenamos
una red neuronal que segmenta, en pares de imágenes Sentinel-2 tomadas antes y después, la superficie desmontada
para alojar pozos, y medimos con el padrón oficial de pozos cuántos meses antes del inicio de la perforación
la detecta.

\subsection{Problema y motivación}
% Enfoque comercial (ADR-0002). Sourcewater va como motivación, no como bibliografía (nota al pie).
Vaca Muerta es una formación geológica de la cuenca Neuquina de la que se extraen petróleo y gas no
convencionales, y su actividad no deja de crecer: según el padrón de pozos de la Secretaría de Energía~\cite{padron},
se pasó de unos 20 pozos perforados por mes en 2018--2019 a unos 36 en 2024--2025. Antes de perforar, la
operadora desmonta un rectángulo de tierra, la \emph{locación}, que aloja varios pozos (en la zona de Añelo, la
mediana es de cuatro). Como la locación se prepara semanas o meses antes de que llegue el equipo de perforación
y es visible desde el espacio, detectarla temprano permite anticipar la actividad antes de que aparezca en los
registros oficiales, que ubican cada pozo como un punto y con demora.

Esa anticipación es inteligencia de mercado: cada empresa conoce su propia actividad, pero no la de la
competencia. A las empresas de servicios (fractura, arena, agua) les permite prever la demanda y salir a buscar
clientes, y a las operadoras, seguir lo que hacen sus vecinas; también les sirve a inversores, analistas y
transportistas para anticipar la producción. En la cuenca Permian de Estados Unidos ya existen servicios
comerciales que venden este tipo de alerta\footnote{Por ejemplo, \emph{DirtWork Alert}, de Sourcewater
(\url{https://www.sourcewater.com}), que informa que las locaciones se desmontan típicamente entre 1 y 4 meses antes
de la perforación. Son cifras de la empresa, sin revisión de pares, por lo que las usamos solo como motivación.}.
El mismo modelo serviría además para el cumplimiento ambiental, por ejemplo para inventariar la huella de la
actividad o verificar la restauración de locaciones abandonadas.

\subsection{Objetivos e hipótesis}
% docs/METODO.md, sección "Hipótesis y experimentos".
El objetivo del trabajo es doble. Por un lado, segmentar en un par de imágenes satelitales las locaciones
que aparecieron entre la primera fecha (T1) y la segunda (T2). Por otro, usar ese modelo sobre las imágenes
mensuales de 2019 a 2023 para medir, con el padrón oficial como referencia, con cuántos meses de anticipación
respecto del inicio de la perforación detecta cada locación nueva. Planteamos cuatro hipótesis, cada una con el
experimento que la pone a prueba:
\begin{itemize}[leftmargin=*,itemsep=1pt]
  \item \textbf{H1 (bandas).} Agregar información infrarroja a RGB mejora la detección, porque ayuda a
  distinguir el suelo desmontado de los cambios naturales de la estepa. Comparamos la red siamesa con tres
  entradas: RGB, seis bandas (RGB, NIR y dos SWIR) y seis bandas más los índices NDVI, NDBI y BSI.
  \item \textbf{H2 (transferencia de aprendizaje).} Con pocos datos, partir de un encoder preentrenado rinde
  más que entrenar desde cero, y ajustarlo entero (\emph{fine-tuning}) más que congelarlo. Entrenamos la red
  siamesa de seis bandas en los tres modos.
  \item \textbf{H3 (generalización).} El desempeño cae en una zona con otras operadoras y otro paisaje.
  Comparamos el modelo final en dos conjuntos de test: bloques de la misma zona de entrenamiento
  (\texttt{test\_anelo}) y una zona externa que el modelo nunca vio (\texttt{test\_externo}).
  \item \textbf{H4 (arquitectura).} Procesar cada fecha con el mismo encoder y comparar sus características
  (red siamesa) funciona mejor que apilar las dos fechas en la entrada (fusión temprana). Comparamos las dos
  arquitecturas con la misma entrada de seis bandas y en el modo \emph{fine-tuning}.
\end{itemize}

\subsection{Estado del arte}
% Con referato. Ramachandran et al. (2024) es la referencia metodológica principal.
\pend{Ramachandran et al.~\cite{ramachandran2024}: locaciones y tanques en Permian y Denver-Julesburg, generalización a regiones nuevas.
U-Net~\cite{ronneberger2015}. Sumar 1--2 papers con referato de segmentación o detección de cambios sobre Sentinel-2.
Modelo de ESRI para el Permian como antecedente (sin acceso).}

\subsection{Imágenes Sentinel-2}
% Sofi. ADR-0004, docs/METODO.md "Bandas e índices".
\pend{Resolución de 10 y 20~m, revisita de 5 días, bandas usadas (B2, B3, B4, B8, B11, B12) y por qué NIR y SWIR
distinguen suelo desmontado de vegetación. Índices NDVI, NDBI y BSI con sus fórmulas.}
\begin{equation}
\mathrm{NDVI} = \frac{B8 - B4}{B8 + B4}
\label{eq:ndvi}
\end{equation}

\subsection{Segmentación semántica y U-Net}
% Agus (marco teórico del modelo).
\pend{Segmentación semántica como clasificación por píxel; arquitectura encoder-decoder con skip connections; encoder ResNet34 preentrenado en ImageNet.}

\subsection{Detección de cambios bitemporal}
% Agus. ADR-0011.
\pend{Dos fechas T1 y T2; fusión temprana (apilar canales) frente a red siamesa (encoder compartido y fusión de características por diferencia).}

\subsection{Funciones de pérdida}
% Agus. resultados/kfold/perdidas.md.
\pend{Desbalance de clases (la locación es $\sim$0,3\,\% de los píxeles); entropía cruzada con pesos, Dice y Lovász.}

% ------------------------------------------------------------------ 2

\section{Metodología}

\pend{Párrafo con el pipeline completo: imágenes mensuales $\rightarrow$ pares $\rightarrow$ etiquetado $\rightarrow$ modelo $\rightarrow$ evaluación $\rightarrow$ backtesting.}

\begin{figure}[t]
\centering
\figpend{Figura: diagrama del pipeline}
\caption{\pend{Pipeline del proyecto.}}
\label{fig:pipeline}
\end{figure}

\subsection{Zonas de estudio y padrón de pozos}
% Sofi. docs/DATASET.md, ADR-0006, data/zonas_estudio.json.
\pend{Padrón de pozos de la Secretaría de Energía (Capítulo IV, CC-BY-4.0); agrupación de pozos en locaciones a 60~m;
locación nueva (primer pozo entre julio de 2019 y diciembre de 2023); Zona Añelo (4$\times$4 bloques) y zona externa;
splits por bloque (train, val, test\_anelo, test\_externo) para evitar fuga espacial.}

\begin{figure}[t]
\centering
\figpend{Figura: mapa de las dos zonas con los bloques coloreados por split y las locaciones nuevas}
\caption{\pend{Zonas de estudio, bloques y splits.}}
\label{fig:zonas}
\end{figure}

\subsection{Imágenes mensuales}
% Sofi. ADR-0010, src/vmseg/gee.py y exportar.py.
\pend{Google Earth Engine; mediana mensual de píxeles despejados con Cloud Score+; enero de 2019 a diciembre de 2023;
meses nublados o con pocas observaciones despejadas.}

\subsection{Pares y etiquetado}
% Sofi. docs/METODO.md "Encuadres y pares", docs/VALIDACION.md.
\pend{Encuadres de 256$\times$256~px; positivos con T1 = mes de perforación $-6$ y T2 = $+1$; negativos con T2 = T1 + 7;
máscaras propuestas con SAM sobre T2 y caída de NDVI; revisión rápida y corrección en QGIS; correcciones de T1 a $-12$ o $-24$
meses y locaciones preexistentes; positivos de otras formaciones en train y val.}

\subsection{Modelos}
% Agus. src/vmseg/siamesa.py, docs/METODO.md (tabla de los 6 modelos).
\pend{Siamesa y fusión temprana con U-Net de smp; entradas RGB, 6 bandas y 6 bandas + índices; modos desde cero,
congelado y fine-tuning; adaptación de la primera capa para más de 3 canales. Solo la clase locación (num\_classes = 2).}

\begin{table}[t]
\centering
\caption{\pend{Los seis modelos entrenados.}}
\label{tab:modelos}
\small
\begin{tabular}{@{}clll@{}}
\toprule
\# & Arquitectura & Entrada & Modo \\
\midrule
1 & Siamesa & 6 bandas & Fine-tuning \\
2 & Fusión temprana & 6 + 6 bandas & Fine-tuning \\
3 & Siamesa & RGB & Fine-tuning \\
4 & Siamesa & 6 bandas + índices & Fine-tuning \\
5 & Siamesa & 6 bandas & Desde cero \\
6 & Siamesa & 6 bandas & Congelado \\
\bottomrule
\end{tabular}
\end{table}

\subsection{Entrenamiento}
% Agus. src/vmseg/entrenar.py, kfold.py, notebooks/entrenar_kfold.ipynb. Devolución del profe (#27).
\pend{K-fold por bloques (5 folds sobre train + val); AdamW, lr, ReduceLROnPlateau, batch, épocas; aumento de datos
(giros y rotaciones de 90°, iguales para T1, T2 y la máscara, sin cambios de brillo ni color); estandarización por canal sobre train;
entrenamiento final con todos los pares y la mediana de las mejores épocas.}

\subsection{Evaluación}
% Bian. src/vmseg/metricas.py.
\pend{Métricas por píxel (IoU de locación, precisión, recall, matriz de confusión) y por objeto; evaluación en test\_anelo y test\_externo.}

\subsection{Backtesting}
% Bian. ADR-0003, ADR-0007, GLOSSARY.md.
\pend{Pares ($m-6$, $m$) sobre las imágenes mensuales; locación capturada (centro dentro de un objeto predicho, con 50~m de margen);
primera detección (dos meses seguidos); anticipación respecto del inicio de perforación; locación detectada a tiempo.}
\begin{equation}
\text{anticipación} = \text{mes de inicio de perforación} - \text{mes de primera detección}
\label{eq:anticipacion}
\end{equation}

% ------------------------------------------------------------------ 3

\section{Resultados y Discusión}

\subsection{Dataset}
% Sofi. data/indice_pares.csv, data/correcciones_pares.csv. Acuerdo: Bian (#19).
\pend{Pares por split (tabla); correcciones de la validación (cuántos a $-12$, $-24$ y preexistentes); acuerdo entre etiquetadoras (IoU, #19).}

\begin{table}[t]
\centering
\caption{\pend{Pares validados por split.}}
\label{tab:dataset}
\small
\begin{tabular}{@{}lrrr@{}}
\toprule
Split & Positivos & Negativos & Total \\
\midrule
train & \pend{75} & \pend{54} & \pend{129} \\
val & \pend{16} & \pend{14} & \pend{30} \\
test\_anelo & \pend{28} & \pend{24} & \pend{52} \\
test\_externo & \pend{45} & \pend{38} & \pend{83} \\
\midrule
Total & \pend{164} & \pend{130} & \pend{294} \\
\bottomrule
\end{tabular}
\end{table}

\begin{figure}[t]
\centering
\figpend{Figura: ejemplos de pares (T1, T2, máscara) con y sin locación}
\caption{\pend{Ejemplos del dataset.}}
\label{fig:ejemplos_dataset}
\end{figure}

\subsection{Elección de la pérdida}
% Agus. resultados/kfold/perdidas.md.
\pend{Tabla de las cuatro pérdidas con IoU de val (media $\pm$ desvío de los 5 folds) y la elegida.}

\subsection{Comparación de modelos (H1, H2 y H4)}
% Agus (#20). Curvas de entrenamiento: loss por iteración y métrica en train y val por época (lo pide la consigna).
\pend{Tabla de los 6 modelos con IoU de locación (media $\pm$ desvío); discusión de cada hipótesis; curvas del mejor modelo.}

\begin{figure}[t]
\centering
\figpend{Figura: curvas de entrenamiento del mejor modelo}
\caption{\pend{Curvas de entrenamiento.}}
\label{fig:curvas}
\end{figure}

\subsection{Evaluación en test (H3)}
% Bian (#21).
\pend{Métricas en test\_anelo y test\_externo; matrices de confusión; ejemplos correctos e incorrectos; caída entre zonas.}

\subsection{Análisis de errores y casos patológicos}
% Agus (#25).
\pend{Falsos positivos (cambios naturales, caminos, otras obras) y falsos negativos (locaciones chicas, meses nublados); hipótesis sobre el comportamiento del modelo.}

\subsection{Backtesting}
% Bian (#26). El número estrella del TP.
\pend{Por zona: \% de locaciones detectadas a tiempo, tarde y nunca; histograma y mediana de la anticipación; comparación entre zonas.}

\begin{figure}[t]
\centering
\figpend{Figura: histograma de la anticipación por zona}
\caption{\pend{Distribución de la anticipación.}}
\label{fig:anticipacion}
\end{figure}

\subsection{Detecciones sin registro}
% Sofi (#22). Se recorta si falta tiempo.
\pend{Muestra de 20 detecciones sin registro clasificadas a ojo: locación sin registrar, otra infraestructura, cambio natural o error.}

% ------------------------------------------------------------------ 4

\section{Conclusiones y trabajo futuro}
% Entre las tres. Trabajo futuro según el #27 y la devolución del profe.
\pend{Conclusiones por hipótesis y la respuesta a la pregunta principal (cuántos meses de anticipación da el satélite).
Trabajo futuro: pérdida robusta al ruido de etiquetas, otras cuencas (Alberta), comparación con ESRI,
encoders preentrenados en imágenes satelitales o con DINOv3.}

% ------------------------------------------------------------------ referencias

\begin{thebibliography}{9}

\bibitem{ramachandran2024}
N.~Ramachandran \emph{et al.}, ``Deep learning for detecting and characterizing oil and gas well pads in satellite imagery,''
\emph{Nature Communications}, vol.~15, 2024. doi: 10.1038/s41467-024-50334-9.

\bibitem{ronneberger2015}
O.~Ronneberger, P.~Fischer y T.~Brox, ``U-Net: Convolutional Networks for Biomedical Image Segmentation,''
en \emph{MICCAI}, 2015.

\bibitem{padron}
Secretaría de Energía de la Nación, ``Producción de petróleo y gas por pozo (Capítulo IV),'' datos.gob.ar, licencia CC-BY-4.0.
\pend{fecha de consulta}

\bibitem{osm}
\pend{OpenStreetMap: qué se usó y cita (© colaboradores de OpenStreetMap, ODbL).}

\bibitem{sentinel}
\pend{Copernicus Sentinel-2 y Cloud Score+ (Google Earth Engine).}

\bibitem{sam}
\pend{A.~Kirillov \emph{et al.}, ``Segment Anything,'' ICCV, 2023 (verificar).}

\end{thebibliography}

\end{document}
